%% kapa: exam=1 round=25 year=2014 period=2 form=A subject=appraisal_law subject_ko=부동산관계법규 seq=15 number=15
%% source: Q-Net 감정평가사 기출문제 (cst00309) artlSeq=5101207 file=extracted/1008/5101207/2014년 제25회 감정평가사 1차시험 기출문제/2교시(A형).hwp
%% conversion: kordoc 4.24.51 -> tools/gen1.py (auto); review=proofread-sonnet; flags=none
\begin{question}{15}
국토의 계획 및 이용에 관한 법령상 도시ㆍ군계획시설사업의 비용 부담에 관한 설명으로 옳지 않은 것은?
\begin{choices}
  \item 도시ㆍ군계획시설사업을 지방자치단체가 하는 경우에는 해당 지방자치단체가 그에 관한 비용을 부담함을 원칙으로 한다.
  \item 국토교통부장관은 그가 시행한 도시ㆍ군계획시설사업으로 현저히 이익을 받는 시ㆍ도가 있으면 사업에 소요된 비용의 50퍼센트를 넘지 않는 범위 안에서 그 비용의 일부를 이익을 받는 시ㆍ도에 부담시킬 수 있다.
  \item 행정청이 아닌 자가 시행하는 도시ㆍ군계획시설사업으로 공공시설의 관리자가 현저한 이익을 받았을 때 시행자는 사업에 소요된 비용의 3분의 1을 넘지 않는 범위 안에서 그 비용의 일부를 관리자에게 부담시킬 수 있다.
  \item 행정청이 시행하는 도시ㆍ군계획시설사업에 드는 비용은 소요비용의 50퍼센트 이하의 범위 안에서 국가예산으로 보조하거나 융자할 수 있다.
  \item 행정청이 아닌 자가 시행하는 도시ㆍ군계획시설사업에 드는 비용은 소요비용의 3분의 1 이하의 범위 안에서 국가 또는 지방자치단체가 보조하거나 융자할 수 있다.
\end{choices}
\end{question}
\answer{3}
